\begin{center}
{\Huge \scshape Qiong Liu, Ph.D.} \\ \vspace{4pt}
\small
\faPhone\ 475-280-1364 \quad
\href{mailto:qiongliu.work@gmail.com}{qiongliu.work@gmail.com} \quad
\href{https://www.linkedin.com/in/qiong-l-529b23196}{LinkedIn} \quad
\href{https://scholar.google.com/citations?user=SvnpHCkAAAAJ&hl=en}{Google Scholar}

Authorized to work in the U.S. without sponsorship
\end{center}

\section{Professional Summary}
\textbf{AI Scientist and Biomedical Imaging Researcher} with expertise in \textbf{cardiovascular AI, cardiac PET imaging, respiratory/cardiac motion correction, physiological signal extraction, tracer kinetic modeling, and quantitative imaging analysis}. Experienced in developing deep learning and physics-informed methods for \textbf{cardiac motion compensation, low-count denoising, parametric imaging, and clinically robust quantitative analysis} in dynamic PET. Strong background in translating cardiovascular imaging workflows into AI-driven solutions through close collaboration with cardiologists, physicians, imaging scientists, and product teams.

\section{Core Expertise}
\textbf{Cardiovascular AI \& Quantitative Imaging:} Cardiac PET imaging, tracer kinetic modeling, parametric imaging, physiological signal analysis, ECG-correlated validation, quantitative biomarker extraction \\
\textbf{Motion Correction \& Signal Processing:} Respiratory/cardiac dual gating, frequency-domain analysis, dynamic PET motion compensation, deformation field fusion \\
\textbf{Machine Learning:} Deep learning, self-supervised learning, physics-informed AI, representation learning, diffusion models, transformers \\
\textbf{Medical Imaging:} PET reconstruction, low-count denoising, dynamic PET, image quality assessment, clinical imaging validation \\
\textbf{Programming:} Python, PyTorch, NumPy, SciPy, MATLAB

\section{Experience}
\textbf{Canon Medical Research USA} \hfill Apr 2025 -- Present \\
\textit{AI Scientist / Reconstruction Scientist, Vernon Hills, IL}
\begin{itemize}[leftmargin=*]
\item Developed AI-driven motion correction and reconstruction methods for \textbf{dynamic cardiac PET imaging} under noisy and low-count acquisition conditions.
\item Designed \textbf{automatic cardiac-respiratory dual gating algorithms} to extract physiological motion signals directly from dynamic PET sequences.
\item Developed \textbf{EdgeFuse}, an edge-guided deformation field fusion framework for cardiac PET motion correction to improve myocardial boundary preservation and quantitative accuracy.
\item Built end-to-end pipelines integrating preprocessing, model training, inference, motion correction, denoising, and downstream quantitative assessment.
\item Collaborated with physicians, imaging scientists, and product teams to evaluate image quality, temporal consistency, quantitative robustness, and clinical usability.
\end{itemize}

\textbf{Canon Medical Research USA} \hfill Jun 2023 -- May 2024 \\
\textit{Research Scientist Intern, Vernon Hills, IL}
\begin{itemize}[leftmargin=*]
\item Developed deep learning methods for low-count PET denoising while preserving quantitative fidelity for downstream dynamic and parametric analysis.
\item Designed evaluation frameworks to assess robustness, generalization, image quality, and quantitative consistency across heterogeneous imaging datasets.
\item Presented technical findings to guide algorithm design and translational imaging research.
\end{itemize}

\textbf{Yale University PET Center} \hfill Aug 2020 -- May 2025 \\
\textit{Ph.D. Researcher, New Haven, CT}
\begin{itemize}[leftmargin=*]
\item Developed AI and quantitative imaging methods for \textbf{dynamic cardiac PET analysis, tracer kinetic modeling, parametric imaging, and low-count denoising}.
\item Applied \textbf{Patlak analysis} and kinetics-informed methods to derive image-based cardiovascular biomarkers from noisy dynamic PET data.
\item Performed quantitative cardiac PET analysis to evaluate treatment-related effects in transthyretin cardiac amyloidosis.
\item Built self-supervised and deep image prior-based denoising methods that preserve kinetic modeling accuracy and downstream parametric quantification.
\item Collaborated with cardiologists, nuclear medicine physicians, and clinical researchers to define clinically meaningful imaging endpoints and validate quantitative results.
\end{itemize}

\textbf{United Imaging Healthcare} \hfill Jul 2019 -- Aug 2019 \\
\textit{Student Intern, Wuhan, China}
\begin{itemize}[leftmargin=*]
\item Supported development and validation of engineering systems for clinical applications.
\end{itemize}

\section{Selected Cardiovascular AI Projects}

\textbf{Automatic Cardiac-Respiratory Dual Gating for Dynamic PET}
\begin{itemize}[leftmargin=*]
\item Developed an AI-driven dual gating framework for simultaneous extraction of respiratory and cardiac motion signals directly from dynamic PET data.
\item Combined autoencoder feature learning, PCA, and frequency-domain analysis to separate physiological motion components.
\item Validated extracted cardiac signals against ECG recordings and visually evaluated respiratory motion using high-temporal-resolution PET mini-frame movies.
\end{itemize}

\textbf{EdgeFuse: Edge-Guided Motion Correction for Cardiac PET}
\begin{itemize}[leftmargin=*]
\item Developed a spatially adaptive deep learning framework for cardiac PET motion correction using edge-guided deformation field fusion.
\item Combined deformation fields with different smoothness constraints to preserve myocardial boundaries while maintaining smooth background motion.
\item Improved motion compensation, lesion localization, and quantitative accuracy in cardiac PET imaging.
\end{itemize}

\textbf{Dynamic Cardiac PET Quantification and Kinetic Modeling}
\begin{itemize}[leftmargin=*]
\item Developed quantitative analysis pipelines for dynamic cardiac PET using Patlak-guided and kinetics-informed methods.
\item Evaluated treatment-related effects in transthyretin cardiac amyloidosis using parametric imaging biomarkers.
\item Improved robustness of quantitative cardiovascular imaging under noisy and low-count acquisition conditions.
\end{itemize}

\textbf{Low-Count and Dynamic PET Denoising}
\begin{itemize}[leftmargin=*]
\item Developed personalized, population-based, and self-supervised denoising frameworks for dynamic and low-count PET imaging.
\item Preserved quantitative fidelity and downstream kinetic modeling accuracy during denoising.
\item Designed evaluation protocols emphasizing image quality, temporal consistency, and parametric quantification.
\end{itemize}

\section{Patents}
\textbf{Edge-Guided Deformation Field Fusion for PET Motion Correction} \hfill Patent Pending \\
\textit{Lead Inventor}
\begin{itemize}[leftmargin=*]
\item Developed an AI-driven cardiac PET motion correction framework that adaptively combines deformation fields to preserve myocardial boundaries and improve quantitative accuracy.
\end{itemize}

\textbf{Automatic Data-Driven Cardiac-Respiratory Gating} \hfill Patent Pending \\
\textit{Lead Inventor}
\begin{itemize}[leftmargin=*]
\item Developed an AI-based dual gating framework for simultaneous extraction of respiratory and cardiac physiological motion signals directly from dynamic PET data.
\end{itemize}

\section{Education}
\textbf{Yale University} \hfill Aug 2020 -- May 2025 \\
Ph.D. in Biomedical Engineering \\
Thesis: Improving Image Quality and Quantification Accuracy for Static and Dynamic PET

\textbf{Huazhong University of Science and Technology} \hfill Sep 2016 -- Jun 2020 \\
B.S. in Biomedical Engineering, Minor in Computer Science \\
GPA: 3.94/4.00

\section{Selected Publications}
\begin{itemize}[leftmargin=*]
\item Liu Q., et al., \textit{Parametric cardiac imaging with $^{18}$F-flutemetamol PET to evaluate the impact of tafamidis in patients with transthyretin cardiac amyloidosis}, \textbf{Journal of Nuclear Medicine}, 2026.
\item Liu Q., et al., \textit{Patlak-guided self-supervised learning for dynamic PET denoising}, \textbf{IEEE Transactions on Radiation and Plasma Medical Sciences}, 2026.
\item Liu Q., et al., \textit{Population-based deep image prior for dynamic PET denoising: a data-driven approach to improve parametric quantification}, \textbf{Medical Image Analysis}, 2024.
\item Liu Q., et al., \textit{A personalized deep learning denoising strategy for low-count PET images}, \textbf{Physics in Medicine \& Biology}, 2022.
\item Liu Q., et al., \textit{Dynamic imaging and tracer kinetic modeling of $^{18}$F-flutemetamol PET for ATTR cardiac amyloidosis patients}, \textbf{SNMMI Oral Presentation}, 2023.
\item Liu Q., et al., \textit{Comparison of network structures for deep learning-based cardiac PET scatter correction}, \textbf{IEEE NSS MIC}, 2024.
\end{itemize}

\section{Awards}
\textbf{Society of Nuclear Medicine and Molecular Imaging (SNMMI)}
\begin{itemize}[leftmargin=*]
\item Young Investigator Award, 1st Place, Cardiovascular Council Clinical Science, 2025.
\item Poster Award, 2nd Place, Physics, Instrumentation \& Data Sciences Council, 2024.
\item Young Investigator Award, 2nd Place, Cardiovascular Council Clinical Science, 2023.
\end{itemize}

\textbf{Mitacs}
\begin{itemize}[leftmargin=*]
\item Globalink Research Award, 2019.
\end{itemize}